% !TEX root = chapter-standalone.tex


\section{Actions of dynamical systems and machines}
\label{sec:system-actions}

In \cref{sec:cat-dyn-lti} and \cref{sec:cat-machines} we saw that LTI systems, Mealy machines, and Moore machines are the morphisms of categories (in the case of Moore machines, of a \SY{semicategory}), in which composition models connecting two systems in series.
These categories describe how systems are \emph{specified} -- by matrices, or by dynamics and read-out functions -- and how they are put together.
They do not yet describe what a system \emph{does}: how it turns the inputs that it receives over time into the outputs that it produces.

In this section we describe what systems do, and the language that we use for this is the language of actions.
Recall from \cref{sec:functors-generalize} that a functor from a category $\CatC$ to $\Set$ can be thought of as an action of $\CatC$: it assigns a set to every object of $\CatC$, and to every morphism of $\CatC$ a function between the corresponding sets, in a way which is compatible with composition.
For a category of systems, this is precisely the form that a description of behavior takes.
The objects are the spaces in which inputs and outputs take their values, and to each of them we assign the set of all signals, or of all streams, with values in it.
To each system we assign the function which maps a signal that the system receives as its input to the signal that it produces as its output.
That this assignment is a functor then has a clear meaning: what a series composition of two systems does is what the first system does, followed by what the second system does.

We proceed in two steps.
To produce outputs, a system has to be started in one of its states.
So we first consider systems together with an initial state, and we see that these form categories of their own.
Then we let these categories act: LTI systems on signals in continuous time, and LTI systems, Mealy machines, and Moore machines on streams in discrete time.
Along the way we will see that one and the same category can act in quite different ways, and that systems which are specified differently can nevertheless act in the same way.

\subsection{Systems with initial states}
\label{sec:systems-with-initial-states}
An LTI system or a machine, as we have defined it, describes how the state changes and which output is produced, whatever the current state may be.
It does not say in which state we begin.
This is a feature: one and the same system can be run from any of its states, and, as we saw in \cref{sec:moore-machines}, different initial states lead, in general, to different outputs.
What transforms inputs into outputs in a definite way is therefore not a system on its own, but a system that has been started: a system \emph{together with} a chosen initial state, in the same way in which a differential equation determines a solution only together with an initial condition.
Such pairs are the protagonists of this section.

Our first step is to see that systems with initial states can be composed, just like the systems themselves.
The idea is simple: to start a series composition of two systems, we start each of the two systems.
The state of the composite system consists of a state of the first system and a state of the second one, and so its initial state consists of the two given initial states.
In this way systems with initial states form categories, which we now define, first for LTI systems and then for machines.

\begin{ctdefinition}[Category $\LTIst$]
    \label{def:LTIst}
    The category $\LTIst$ of LTI systems with initial state is the category given by:

    \constit

    \begin{enumerate}
        \item Objects: the real coordinate spaces $\reals^n$, for $n \setin \natnumbers$, as in $\LTI$.
        \item Morphisms: a morphism from $\reals^n$ to $\reals^m$ is a pair $\tup{\mora, \prstart}$, where
              \begin{equation}
                  \mora = \tupp{\reals^k, \mat{A}, \mat{B}, \mat{C}, \mat{D}} \colon \reals^n \mto \reals^m
              \end{equation}
              is a morphism of $\LTI$ (\cref{def:LTICat}), and $\prstart \setin \reals^k$ is an element of its state space, the \emph{initial state}.
        \item Composition: given morphisms $\tup{\mora, \prstart_\mora}$ from $\reals^n$ to $\reals^m$ and $\tup{\morb, \prstart_\morb}$ from $\reals^m$ to $\reals^p$, where the state spaces of $\mora$ and $\morb$ are $\reals^k$ and $\reals^l$, their composition is
              \begin{equation}
                  \label{eq:LTIst-composition}
                  \tup{\mora, \prstart_\mora} \mthen \tup{\morb, \prstart_\morb} \definedas \tup{\morab, \begin{bmatrix} \prstart_\mora \\ \prstart_\morb \end{bmatrix}},
              \end{equation}
              where $\morab$ is the composition in $\LTI$ (\cref{eq:LTI-composition}), and where the vector obtained by stacking $\prstart_\mora$ on top of $\prstart_\morb$ is an element of the state space $\reals^{k+l}$ of $\morab$.
        \item Identity morphisms: the identity morphism of $\reals^n$ is the pair
              \begin{equation}
                  \tup{\catidat{\reals^n}, \mat{0}^{\matdim{0}{1}}},
              \end{equation}
              where $\catidat{\reals^n}$ is the identity morphism of $\reals^n$ in $\LTI$, whose state space is $\reals^0$, and $\mat{0}^{\matdim{0}{1}}$ is the only element of $\reals^0$: the column vector with no entries.
    \end{enumerate}
\end{ctdefinition}

Composition in $\LTIst$ is \SY{associative} and unital because composition in $\LTI$ is, and because stacking column vectors is \SY{associative}, and stacking a column vector with the column vector with no entries does not change it (compare \cref{rem:LTI-state-space-trick}).

For machines we proceed in the same way.
A \emph{Mealy machine with initial state} is a pair $\tup{\mora, \prstart}$ of a Mealy machine $\mora = \tupp{\prin, \prst, \prout, \prdyn, \prreadout}$ (\cref{def:mealy-machine}) and an element $\prstart \setin \prst$ of its set of states, the \emph{initial state}.
Two Mealy machines with initial state $\tup{\mora, \prstart}$ and $\tup{\mora', \prstart'}$, with the same set of inputs and the same set of outputs, are \emph{isomorphic} if there is an isomorphism $\mapc$ of Mealy machines from $\mora$ to $\mora'$ (\cref{sec:cat-machines}) which is compatible with the initial states, that is, such that $\mapc(\prstart) = \prstart'$.
As in \cref{sec:cat-machines}, this is an \SY{equivalence relation}, and we write $\stylemorph{[}\mora, \prstart\stylemorph{]}$ for the equivalence class of $\tup{\mora, \prstart}$.

\begin{ctdefinition}[Category $\MealyModst$]
    \label{def:MealyModst}
    The category $\MealyModst$ of Mealy machines with initial state, up to isomorphism, is the category given by:

    \constit

    \begin{enumerate}
        \item Objects: all sets.
        \item Morphisms: a morphism from a set $\prin$ to a set $\prout$ is an equivalence class $\stylemorph{[}\mora, \prstart\stylemorph{]}$ of Mealy machines with initial state whose set of inputs is $\prin$ and whose set of outputs is $\prout$.
        \item Composition: given morphisms $\stylemorph{[}\mora, \prstart_\mora\stylemorph{]}$ and $\stylemorph{[}\morb, \prstart_\morb\stylemorph{]}$, where the set of outputs of $\mora$ is the set of inputs of $\morb$, their composition is
              \begin{equation}
                  \label{eq:MealyModst-composition}
                  \stylemorph{[}\mora, \prstart_\mora\stylemorph{]} \mthen \stylemorph{[}\morb, \prstart_\morb\stylemorph{]} \definedas \stylemorph{[}\morab, \tup{\prstart_\mora, \prstart_\morb}\stylemorph{]},
              \end{equation}
              where $\morab$ is the Mealy machine defined in \cref{eq:mealy-composition}, whose set of states is $\prst_{\mora} \cartprod \prst_{\morb}$.
        \item Identity morphisms: the identity morphism of a set $\prin$ is $\stylemorph{[}\more, \ast\stylemorph{]}$, where $\more$ is the Mealy machine with the single state $\ast$ which represents the identity morphism of $\prin$ in $\MealyMod$ (\cref{def:MealyCat}).
    \end{enumerate}
\end{ctdefinition}

That \cref{def:MealyModst} does define a category is checked in the same way as for $\MealyMod$ (\cref{rem:MealyMod-category}): the isomorphisms of Mealy machines that are used there -- pairs of isomorphisms of the two machines that are composed, re-bracketing of states, and forgetting the component $\ast$ -- are all compatible with the initial states.

In exactly the same way, Moore machines with initial state, up to isomorphism, form a \SY{semicategory}, which we denote by $\MooreModst$.
Its morphisms are the equivalence classes $\stylemorph{[}\mora, \prstart\stylemorph{]}$ of pairs of a Moore machine $\mora$ (\cref{def:moore-machine-first}) and an element $\prstart$ of its set of states, and its composition is given by the same formula as in \cref{eq:MealyModst-composition}, where now $\morab$ is the Moore machine of \cref{eq:MooreMod-composition}.

In each of the three cases, forgetting the initial state is a functor (for Moore machines, a \SY{semifunctor}) to the category, or semicategory, that we started with: from $\LTIst$ to $\LTI$, from $\MealyModst$ to $\MealyMod$ and from $\MooreModst$ to $\MooreMod$ respectively.

\begin{remark}[Why the initial state is part of the morphism]
    \label{rem:why-initial-states}
    One might ask whether the categories of systems \emph{without} initial states could act on signals and on streams directly.
    For this, we would have to select one initial state for every system, once and for all, and in such a way that the state selected for a composite system consists of the states selected for its two parts.
    For LTI systems this is possible: their state spaces are coordinate spaces, and we can always select the zero vector.
    For machines it is not: a set of states has, in general, no distinguished element, and it may even be empty.
    Treating the initial state as part of the morphism works in all cases, and it lets one and the same system act in many ways, one for each of its states.
\end{remark}


\subsection{LTI actions on signals}

With these categories at hand, we come to the second step, and let systems with initial states act.
We begin with LTI systems, and with signals in continuous time.

\begin{example}[LTI systems act on signals]
    \label{exa:lti-action-signals}
    Recall the category $\LTIst$ (\cref{def:LTIst}): its objects are the real coordinate spaces $\reals^n$, and a morphism from $\reals^n$ to $\reals^m$ is a pair $\tup{\mora, \prstart}$ of an LTI system
    \begin{equation}
        \mora = \tupp{\reals^k, \mat{A}, \mat{B}, \mat{C}, \mat{D}}
    \end{equation}
    with input space $\reals^n$, output space $\reals^m$, and state space $\reals^k$, and an initial state $\prstart \setin \reals^k$.
    We now take the \emph{dynamical systems} point of view on such a tuple (\cref{def:LTI_syst_first_time}): it specifies a system in continuous time, which, once it is started in the state $\prstart$, turns a signal with values in $\reals^n$ into a signal with values in $\reals^m$.
    This defines an action of $\LTIst$, that is, a functor
    \begin{equation}
        \funa \colon \LTIst \fto \Set.
    \end{equation}

    On objects, $\funa$ assigns to each coordinate space the set of signals with values in it, which we model as continuously differentiable functions of time:
    \begin{equation}
        \funaob(\reals^n) \definedas \diffun(\nonNegReals, \reals^n).
    \end{equation}
    On morphisms, $\funa$ assigns to a pair $\tup{\mora, \prstart}$ as above the function
    \begin{equation}
        \funamor(\mora, \prstart) \colon \diffun(\nonNegReals, \reals^n) \sto \diffun(\nonNegReals, \reals^m)
    \end{equation}
    which maps an input signal $\mat{u}$ to the output signal $\mat{y}$ given by
    \begin{equation}
        \label{eq:lti-action-signals-output}
        \mat{y}(t) \definedas \mat{C}\mat{s}(t) + \mat{D}\mat{u}(t),
    \end{equation}
    where $\mat{s} \colon \nonNegReals \sto \reals^k$ is the unique solution of the initial value problem
    \begin{equation}
        \label{eq:lti-action-signals-ivp}
        \begin{cases}
            \dot{\mat{x}}(t) = \mat{A}\mat{x}(t) + \mat{B}\mat{u}(t), \\
            \mat{x}(0) = \prstart.
        \end{cases}
    \end{equation}
    Such a solution always exists and is unique; explicitly,
    \begin{equation}
        \mat{s}(t) = e^{\mat{A} t} \prstart + \int_{0}^{t} e^{\mat{A} (t - \tau)}\mat{B}\mat{u}(\tau) \, d\tau.
    \end{equation}

    Let us see why the functor conditions hold.
    Let $\tup{\mora, \prstart_\mora} \colon \reals^n \mto \reals^m$ and $\tup{\morb, \prstart_\morb} \colon \reals^m \mto \reals^p$ be LTI systems with initial states, let $\mat{u}$ be an input signal for $\mora$, and let $\mat{y}_\mora \definedas \funamor(\mora, \prstart_\mora)(\mat{u})$.
    Let $\mat{s}_\mora$ be the solution of the initial value problem of $\mora$, with initial state $\prstart_\mora$, for the input $\mat{u}$, and let $\mat{s}_\morb$ be the solution of the initial value problem of $\morb$, with initial state $\prstart_\morb$, for the input $\mat{y}_\mora$.
    Stacking $\mat{s}_\mora$ on top of $\mat{s}_\morb$, we obtain a function with values in the state space of the composite system $\morab$, and this function solves the initial value problem of the composite system $\morab$, with the stacked initial state of \cref{eq:LTIst-composition}, for the input $\mat{u}$: this is precisely how the block matrices in \cref{eq:LTI-composition} were chosen.
    By uniqueness, it is \emph{the} solution, and so the output signal of $\morab$ is
    \begin{equation}
        \mat{D}_\morb\mat{C}_\mora\mat{s}_\mora(t) + \mat{C}_\morb\mat{s}_\morb(t) + \mat{D}_\morb\mat{D}_\mora\mat{u}(t)
        = \mat{C}_\morb\mat{s}_\morb(t) + \mat{D}_\morb\mat{y}_\mora(t),
    \end{equation}
    which is the output signal of $\morb$ for the input $\mat{y}_\mora$.
    Hence
    \begin{equation}
        \funamor(\tup{\mora, \prstart_\mora} \mthen \tup{\morb, \prstart_\morb}) = \funamor(\mora, \prstart_\mora) \mthenof{\Set} \funamor(\morb, \prstart_\morb).
    \end{equation}
    The identity morphism of $\reals^n$ has the state space $\reals^0$ and $\mat{D} = \mat{1}^{\matdim{n}{n}}$, so \cref{eq:lti-action-signals-output} reduces to $\mat{y}(t) = \mat{u}(t)$, and $\funamor(\catidat{\reals^n})$ is the identity function.
\end{example}

\begin{lemma}
    \label{lem:equivsystequivact}
    Equivalent LTI systems act on signals in the same way, when they are started in states that correspond to each other: if $\mora, \morb \colon \reals^n \mto \reals^m$ are equivalent LTI systems (\cref{def:equivalence_lti}), with equivalence transformation $\mat{T}$, then
    \begin{equation}
        \funamor(\mora, \prstart) = \funamor(\morb, \mat{T}\prstart)
    \end{equation}
    for every initial state $\prstart$ of $\mora$, where $\funa \colon \LTIst \fto \Set$ is the action of \cref{exa:lti-action-signals}.
\end{lemma}

\begin{proof}
    Consider two equivalent LTI systems
    \begin{equation}
        \mora = \ltituple{\mora},\quad \morb = \ltituple{\morb},
    \end{equation}
    with equivalence transformation $\mat{T}$, let $\prstart$ be an initial state of $\mora$, and let $\mat{u}$ be an input signal.
    The initial value problem that $\mora$ poses, for the initial state $\prstart$, reads:
    \begin{equation}
        \label{eq:first_lti_equiv}
        \begin{aligned}
            \dot{\mat{x}}(t) & =\mat{A}_{\mora}\mat{x}(t)+\mat{B}_{\mora}\mat{u}(t), \\
            \mat{x}(0)       & = \prstart.
        \end{aligned}
    \end{equation}
    Let~$\mat{s}$ be the solution of \cref{eq:first_lti_equiv}, and consider~$\mat{r}(t) \definedas \mat{T}\mat{s}(t)$.
    Then, we have
    \begin{equation}
        \begin{aligned}
            \dot{\mat{r}}(t) & =\mat{T}\dot{\mat{s}}(t) \\
                             & =\mat{T}\mat{A}_\mora \mat{s}(t)+\mat{T}\mat{B}_\mora \mat{u}(t) \\
                             & =\mat{T}\mat{A}_\mora \mat{T}^{-1}\mat{r}(t)+\mat{T}\mat{B}_\mora \mat{u}(t) \\
                             & =\mat{A}_\morb\mat{r}(t)+\mat{B}_\morb \mat{u}(t),
        \end{aligned}
    \end{equation}
    and $\mat{r}(0) = \mat{T} \prstart$.
    So $\mat{r}$ is the unique solution of the initial value problem that $\morb$ poses, for the initial state $\mat{T}\prstart$ and the same input signal $\mat{u}$.
    Therefore, the action of the system $\morb$ with initial state $\mat{T}\prstart$ is:
    \begin{equation}
        \begin{aligned}
            \funamor(\morb, \mat{T}\prstart)(\mat{u})(t) & =\mat{C}_\morb \mat{r}(t)+ \mat{D}_\morb \mat{u}(t) \\
                                        & =\mat{C}_\mora \mat{T}^{-1}\mat{T}\mat{s}(t)+\mat{D}_\mora \mat{u}(t) \\
                                        & =\mat{C}_\mora \mat{s}(t)+\mat{D}_\mora \mat{u}(t) \\
                                        & =\funamor(\mora, \prstart)(\mat{u})(t).
        \end{aligned}
    \end{equation}
\end{proof}

\begin{remark}
    \label{rem:same-action-not-equivalent}
    The converse of \cref{lem:equivsystequivact} does not hold: two LTI systems which, for suitable initial states, have the same action on signals are not necessarily equivalent.
    For a simple counterexample, consider the LTI system with initial state $\prstart = \mat{0}$ and
    \begin{equation}
        \begin{aligned}
            \dot{\mat{x}}(t) & =\begin{bmatrix}
                                    1 & 0 \\ 0& 1
                                \end{bmatrix} \mat{x}(t)+\begin{bmatrix}
                                                             1 \\ 0
                                                         \end{bmatrix} \mat{u}(t) \\
            \mat{y}(t)       & =\begin{bmatrix}
                                    1 & 0
                                \end{bmatrix} \mat{x}(t)+\mat{u}(t).
        \end{aligned}
    \end{equation}
    Its action on signals is the same as that of any system, again with initial state $\mat{0}$, of the form
    \begin{equation}
        \begin{aligned}
            \dot{\mat{x}}(t) & =\begin{bmatrix}
                                    1 & 0 \\ 0& \alpha
                                \end{bmatrix} \mat{x}(t)+\begin{bmatrix}
                                                             1 \\ 0
                                                         \end{bmatrix} \mat{u}(t) \\
            \mat{y}(t)       & =\begin{bmatrix}
                                    1 & 0
                                \end{bmatrix} \mat{x}(t)+\mat{u}(t),
        \end{aligned}
    \end{equation}
    where~$\alpha \setin \reals$, because the output~$\mat{y}$ does not depend on the second component of~$\mat{x}$.
    However, for $\alpha \neq 1$ there is no invertible matrix $\mat{T}$ which relates these systems: the first one has $\mat{A} = \idmat$, and $\mat{T}\idmat\mat{T}^{-1} = \idmat$ for every invertible $\mat{T}$.
    We note that control theory has developed tools to deal with hidden parts of the dynamics.
    Examples are the notion of controllability (ability of an actuator/input to control all states of a system) and observability (ability to estimate all states of the system).
\end{remark}

\subsection{LTI actions on streams}

The tuples that make up the morphisms of $\LTIst$ do not, by themselves, know anything about continuous time.
We can just as well take the \emph{machines} point of view on them, and this gives a second action of the very same category.

\begin{example}[LTI systems act on streams]
    \label{exa:lti-action-streams}
    A \emph{stream} in a set $\setA$ is an infinite list $\tup{a_0, a_1, a_2, \dots}$ of elements of $\setA$, that is, a function $\natnumbers \sto \setA$.
    We write $\setA^{\natnumbers}$ for the set of streams in $\setA$.

    We define a functor
    \begin{equation}
        \funb \colon \LTIst \fto \Set
    \end{equation}
    as follows.
    On objects, $\funb$ assigns to each coordinate space the set of streams of vectors in it:
    \begin{equation}
        \funbob(\reals^n) \definedas (\reals^n)^{\natnumbers}.
    \end{equation}
    On morphisms, $\funb$ assigns to an LTI system $\mora = \tupp{\reals^k, \mat{A}, \mat{B}, \mat{C}, \mat{D}} \colon \reals^n \mto \reals^m$ with initial state $\prstart \setin \reals^k$ the function
    \begin{equation}
        \funbmor(\mora, \prstart) \colon (\reals^n)^{\natnumbers} \sto (\reals^m)^{\natnumbers}
    \end{equation}
    which maps a stream $\mat{u} = \tup{\mat{u}_0, \mat{u}_1, \mat{u}_2, \dots}$ of inputs to the stream $\mat{y} = \tup{\mat{y}_0, \mat{y}_1, \mat{y}_2, \dots}$ of outputs defined, step by step, by
    \begin{equation}
        \label{eq:lti-action-streams}
        \mat{x}_0 \definedas \prstart,
        \qquad
        \mat{x}_{i+1} \definedas \mat{A}\mat{x}_i + \mat{B}\mat{u}_i,
        \qquad
        \mat{y}_i \definedas \mat{C}\mat{x}_i + \mat{D}\mat{u}_i.
    \end{equation}
    In other words, we now read the tuple $\mora$ as a discrete-time LTI system (\cref{def:LTI-discrete}), or, equivalently, as a Mealy machine with set of inputs $\reals^n$, set of states $\reals^k$, and set of outputs $\reals^m$ (\cref{exa:discrete-lti-mealy}), and we let this machine run forever, starting in the state $\prstart$.
    No differential equation needs to be solved: the matrix $\mat{A}$ now describes the next state, and not the rate of change of the state.

    The functor conditions are checked as in \cref{exa:lti-action-signals}, and more easily.
    If we feed the outputs of $\mora$ into $\morb$ at every step, then the stacked states satisfy the recursion \cref{eq:lti-action-streams} for the block matrices in \cref{eq:LTI-composition} and the stacked initial state of \cref{eq:LTIst-composition}; the computation is the same as in the continuous-time case, with $\mat{x}_{i+1}$ in place of $\dot{\mat{x}}(t)$.
    Hence
    \begin{equation}
        \funbmor(\tup{\mora, \prstart_\mora} \mthen \tup{\morb, \prstart_\morb}) = \funbmor(\mora, \prstart_\mora) \mthenof{\Set} \funbmor(\morb, \prstart_\morb).
    \end{equation}
    The identity morphism of $\reals^n$ again yields $\mat{y}_i = \mat{u}_i$, the identity function on streams.
\end{example}

The two examples have the same source category, but they are different functors: the category $\LTIst$ only records how systems are specified, in which state they start, and how they are connected in series, while each of the two actions gives these specifications a meaning.
The same morphism can mean quite different things under the two actions.
Consider, for instance, the LTI system $\tupp{\reals^1, 0, 1, 1, 0} \colon \reals^1 \mto \reals^1$, in which all four matrices have a single entry, together with the initial state $0$.
Under the action $\funa$ on signals, it is an integrator: it maps a signal $u$ to the signal $t \mapsto \int_0^t u(\tau) \, d\tau$.
Under the action $\funb$ on streams, it is a delay by one step: it maps a stream $\tup{u_0, u_1, u_2, \dots}$ to the stream $\tup{0, u_0, u_1, \dots}$.

\subsection{Mealy machine actions}

In \cref{exa:lti-action-streams}, linearity played no role.
The same construction works for arbitrary Mealy machines.

\begin{example}[Mealy machines act on streams]
    \label{exa:mealy-action-streams}
    Recall from \cref{def:MealyModst} the category $\MealyModst$: its objects are sets, and a morphism from a set $\prin$ to a set $\prout$ is an equivalence class $\stylemorph{[}\mora, \prstart\stylemorph{]}$ of Mealy machines $\mora = \tupp{\prin, \prst, \prout, \prdyn, \prreadout}$ with initial state $\prstart \setin \prst$, up to isomorphism.
    There is a functor
    \begin{equation}
        \func \colon \MealyModst \fto \Set
    \end{equation}
    defined as follows.
    On objects, $\func$ assigns to each set the set of streams in it:
    \begin{equation}
        \funobspace{H}(\prin) \definedas \prin^{\natnumbers}.
    \end{equation}
    On morphisms, $\func$ assigns to $\stylemorph{[}\mora, \prstart\stylemorph{]}$ the function
    \begin{equation}
        \funmorspace{H}(\stylemorph{[}\mora, \prstart\stylemorph{]}) \colon \prin^{\natnumbers} \sto \prout^{\natnumbers}
    \end{equation}
    which maps a stream $\tup{\prineln{0}, \prineln{1}, \prineln{2}, \dots}$ of inputs to the stream $\tup{\prouteln{0}, \prouteln{1}, \prouteln{2}, \dots}$ of outputs defined, step by step, by
    \begin{equation}
        \label{eq:mealy-action-streams}
        \prsteln{0} \definedas \prstart,
        \qquad
        \prsteln{i+1} \definedas \prdyn(\prineln{i}, \prsteln{i}),
        \qquad
        \prouteln{i} \definedas \prreadout(\prineln{i}, \prsteln{i}).
    \end{equation}
    This is the stream version of \cref{eq:mealy-list-output}, with the entries numbered starting from $0$.

    Since a morphism is an equivalence class, we need to check that the result does not depend on the chosen representative $\tup{\mora, \prstart}$.
    Let $\mapc$ be an isomorphism from $\tup{\mora, \prstart}$ to a Mealy machine with initial state $\tup{\mora', \prstart'}$, and let $\prsteln{i}'$ and $\prouteln{i}'$ be the states and outputs of $\mora'$, with initial state $\prstart'$, for the same stream of inputs.
    By induction on $i$, the compatibility of $\mapc$ with the initial states and the dynamics gives $\prsteln{i}' = \mapc(\prsteln{i})$, and then the compatibility with the read-outs gives $\prouteln{i}' = \prouteln{i}$ (compare \cref{sec:homomorphisms-mealy}).

    For the functor conditions, consider a composite as in \cref{eq:MealyModst-composition}, represented by the machine $\morab$ of \cref{eq:mealy-composition} with the initial state $\tup{\prstart_\mora, \prstart_\morb}$.
    Its state at step $i$ is the pair consisting of the state of $\mora$ at step $i$ for the given inputs, and the state of $\morb$ at step $i$ when $\morb$ receives the outputs of $\mora$ as its inputs; this follows by induction from the formula for $\prdyn_{\morab}$.
    The formula for $\prreadout_{\morab}$ then says that the $i$-th output of $\morab$ is the $i$-th output of $\morb$.
    Hence
    \begin{equation}
        \funmorspace{H}(\stylemorph{[}\mora, \prstart_\mora\stylemorph{]} \mthen \stylemorph{[}\morb, \prstart_\morb\stylemorph{]}) = \funmorspace{H}(\stylemorph{[}\mora, \prstart_\mora\stylemorph{]}) \mthenof{\Set} \funmorspace{H}(\stylemorph{[}\morb, \prstart_\morb\stylemorph{]}).
    \end{equation}
    The identity morphism of $\prin$ is represented by the machine which passes each input straight through to the output, so it acts as the identity function on $\prin^{\natnumbers}$.
\end{example}

\begin{remark}
    \label{rem:lti-streams-via-mealy}
    The action of $\LTIst$ on streams in \cref{exa:lti-action-streams} is a special case of \cref{exa:mealy-action-streams}: reading an LTI system $\mora \colon \reals^n \mto \reals^m$ with initial state $\prstart$ as a Mealy machine with the same initial state, as we did there, and then applying $\funmorspace{H}$ to its equivalence class gives precisely $\funbmor(\mora, \prstart)$.

    In fact, this reading is itself a functor $\LTIst \fto \MealyModst$: it is defined in the same way as the functor $\LTI \fto \MealyMod$ of \cref{exa:lti-to-mealy-functor}, and it keeps the initial state.
    Once we know how to compose functors (\cref{def:functor-composition}), we will be able to say that the action $\funb$ is the composition of that functor with the action $\func$.
\end{remark}


\subsection{Moore machine actions.}
\label{sec:Moore-semicat-actions}

Recall that Moore machines only form a semicategory; however they act on streams in a manner analogous to the case of Mealy machines treated above. 


\begin{example}[Moore machines act on streams]
    \label{exa:moore-action-streams}
    Recall the semicategory $\MooreModst$ (\cref{sec:systems-with-initial-states}): its objects are sets, and a morphism from a set $\prin$ to a set $\prout$ is an equivalence class $\stylemorph{[}\mora, \prstart\stylemorph{]}$ of Moore machines $\mora = \tupp{\prin, \prst, \prout, \prdyn, \prreadout}$ with initial state $\prstart \setin \prst$, up to isomorphism.
    There is a semifunctor
    \begin{equation}
        \funa \colon \MooreModst \fto \Set
    \end{equation}
    defined as follows.
    On objects, $\funa$ assigns to each set the set of streams in it (\cref{exa:lti-action-streams}):
    \begin{equation}
        \funaob(\prin) \definedas \prin^{\natnumbers}.
    \end{equation}
    On morphisms, $\funa$ assigns to $\stylemorph{[}\mora, \prstart\stylemorph{]}$ the function
    \begin{equation}
        \funamor(\stylemorph{[}\mora, \prstart\stylemorph{]}) \colon \prin^{\natnumbers} \sto \prout^{\natnumbers}
    \end{equation}
    which maps a stream $\tup{\prineln{0}, \prineln{1}, \prineln{2}, \dots}$ of inputs to the stream $\tup{\prouteln{0}, \prouteln{1}, \prouteln{2}, \dots}$ of outputs defined, step by step, by
    \begin{equation}
        \label{eq:moore-action-streams}
        \prsteln{0} \definedas \prstart,
        \qquad
        \prsteln{i+1} \definedas \prdyn(\prineln{i}, \prsteln{i}),
        \qquad
        \prouteln{i} \definedas \prreadout(\prsteln{i}).
    \end{equation}
    This is the stream version of \cref{eq:moore-list-output}.
    As in \cref{exa:mealy-action-streams}, the result does not depend on the chosen representative $\tup{\mora, \prstart}$, since isomorphic Moore machines, started in initial states which correspond to each other, produce the same outputs.

    For the semifunctor condition, consider a composite, represented by the machine $\morab$ of \cref{eq:MooreMod-composition} with the initial state $\tup{\prstart_\mora, \prstart_\morb}$.
    Its state at step $i$ is the pair consisting of the state of $\mora$ at step $i$ for the given inputs, and the state of $\morb$ at step $i$ when $\morb$ receives the outputs of $\mora$ as its inputs; this follows by induction from the formula \cref{eq:moore-comp-naive-2} for $\prdyn_{\morab}$.
    Since the read-out of $\morab$ is the read-out of $\morb$, the $i$-th output of $\morab$ is the $i$-th output of $\morb$.
    Hence
    \begin{equation}
        \funamor(\stylemorph{[}\mora, \prstart_\mora\stylemorph{]} \mthen \stylemorph{[}\morb, \prstart_\morb\stylemorph{]}) = \funamor(\stylemorph{[}\mora, \prstart_\mora\stylemorph{]}) \mthenof{\Set} \funamor(\stylemorph{[}\morb, \prstart_\morb\stylemorph{]}).
    \end{equation}

    Note that the first output, $\prouteln{0} = \prreadout(\prstart)$, does not depend on the inputs at all.
    So, as soon as $\prin$ has more than one element, no Moore machine acts as the identity function on $\prin^{\natnumbers}$; this matches the fact that $\MooreMod$ has no identity morphisms (\cref{rem:MooreMod-no-identities}), and neither does $\MooreModst$.
\end{example}


